\section{All projective tests on the small-r weak-C2 source opening}
\label{sec:r11-small-r-remote}
\begin{theorem}
Retain the original physical weak-C2 source, exact stationary contact,
all strict source gaps, $n\ge0$, and $d=1$. Suppose $B,z,y,r>0$
are infinitesimal, $G/B,j,y/B,L/B$ are positive bounded units,
and $H/y$ is infinitesimal. Then the entire original parent is
finite tensor SOS at every independent projective testing pair,
with all signs, literal endpoints, and real or complex states.
There is no comparison of independent source and testing rates.
\end{theorem}
Write $a=z/B$, $t=y/B$, $h=H/B$. Positivity of
$G/B=1-a(1-r-H)$ makes $a$ bounded; its unit gap gives
$0\le\rho=\operatorname{st}(a)<1$. Put
$\tau=\operatorname{st}(t)>1$. The preceding local lemma proves
the source estimates, the exact full-parent five-root identity,
and its strict residue for every bounded testing pair. In particular
\[
 \mu_0=\frac1{2(1-\rho^2)},\quad m_0=\rho^2\mu_0,\quad
 C_0=\frac{\tau-\rho}{2(1+\rho)},\quad
 D_0=\frac1{2(1+\rho)(\tau-1)}.                 \tag{SR1}
\]
Here $\mu=m/a^2$, while $n,A$ are unbounded, independently.
This section supplies the complete remote join; no remote theorem
requiring bounded $n$ is imported as a domain admission.

\paragraph{Two normalized hard state coordinates.}
Use $\widehat X=\sqrt n X$, $\widehat Y=\sqrt A Y$, with
$e_x=n^{-1/2}$ and $e_y=A^{-1/2}$ positive infinitesimal.
Substitute $X=e_x\widehat X$, $Y=e_y\widehat Y$ into the
original full parent. Its $nX^2$ and $AY^2$ packets become the
exact hard squares, and every remaining coefficient is bounded.
Every cross term with these two normalized states has zero residue.
At bounded test residues $u_0,v_0$, put $U=u_0^2$, $V=v_0^2$.
The entire residue of $F/4$ is
\[
 (1+U)(1+V)(\widehat X^2+\widehat Y^2)+\Psi(W,Z),
\]
where
\[
\begin{split}
 \Psi={}&m_0(v_0W-u_0Z)^2+C_0(1+V)U W^2
 +(D_0+\tfrac12)(u_0W-v_0Z)^2\\
 &+\tfrac12(1+V)UZ^2-\tfrac12(1+U)VZ^2.
\end{split}                                                   \tag{SR2}
\]
This retains the complete signed core. If $P_2$ is its matrix on
$(W,Z)$, direct expansion using (SR1) gives
\[
\begin{split}
 \det P_2={}&D_0\mu_0(\rho V-\tau U)^2\\
 &+\frac{C_0(2m_0+1)}2U^2V+C_0D_0UV^2.          
\end{split}                                                   \tag{SR3}
\]
All three displayed coefficients are strictly positive.
For $U>0$, this determinant is positive for every $V\ge0$:
if $V=0$ its first term is positive; if $V>0$ its second is.
The $W$ diagonal $m_0V+[C_0(1+V)+D_0+1/2]U$ is positive.
Thus the full four-state residue is strictly positive.
For $U=0,V>0$ the same holds when $\rho>0$; here the two
remaining pivots are $m_0V,D_0V$.
Expanding around the real residues and applying the finite separated
lift (WA10) proves every such bounded testing chart.
All bounded pairs are already covered by the preceding strengthened
lemma. The determinant gives a second direct check away from its
vanishing locus.

\paragraph{The joint small-a and small-r corner.}
Only $\rho=0$, $u$ infinitesimal or zero, and nonzero or infinite
$v$ remains outside those arguments. The complete identities
(WE8) and (WE10) apply algebraically at arbitrary $n,A$.
In those exact identities make the same pure-source substitutions
$X=e_x\widehat X$, $Y=e_y\widehat Y$, using the five roots
\[
 (S,R,\widehat X,\widehat Y,Z)=(aW,uW,\sqrt n X,\sqrt A Y,Z).
\]
All resulting coefficients are bounded polynomials in the testing
coordinates with pure-source coefficients. The factors $nX^2$
and $AY^2$ again become the exact normalized squares; no term is
removed. For $v_0\ne0$ finite, the separate pivots on
$(S,\widehat X,\widehat Y)$ are
\[
 \mu_0V,\qquad 1+V,\qquad 1+V.
\]
The remaining block on $(R,Z)$ is
\[
 \begin{pmatrix}
 C_0(1+V)+D_0+1/2&-(D_0+1/2)v_0\\
 -(D_0+1/2)v_0&D_0V
 \end{pmatrix},\qquad \det=C_0D_0V^2>0.          \tag{SR4}
\]
Indeed at $\rho=0$, $C_0=\tau/2$,
$D_0=1/[2(\tau-1)]$, and
$C_0D_0-D_0/2-1/4=0$.
For $v$ unbounded put $q=1/v$ and use the polynomial parent
$q^2F(u,1/q)$. At $u=q=0$ the full five-root residue of the
normalized (WE10) is diagonal with pivots
\[
 (\mu_0,C_0,1,1,D_0)>0.                           \tag{SR5}
\]
The same finite separated lift covers both charts, including
literal $u=0$ and literal $q=0$. The roots may be dependent;
a strict Gram identity remains a valid finite SOS on them.

\paragraph{Every other projective endpoint.}
For $u$ unbounded use $p=1/u$ and the exact polynomial
$p^2F(1/p,v)$. With $v$ bounded its normalized full residue is
\[
 [C_0(1+V)+D_0+1/2]W^2
 +(1+V)(\widehat X^2+\widehat Y^2)+(m_0+1/2)Z^2.
                                                               \tag{SR6}
\]
It is strictly positive for every $V\ge0$, whether $\rho$ vanishes
or not. For $v$ unbounded and $u$ bounded the residue is
\[
 (m_0+C_0U)W^2+(1+U)(\widehat X^2+\widehat Y^2)+D_0Z^2.
                                                               \tag{SR7}
\]
This is strictly positive if $\rho>0$ or $U>0$; the missing
case is precisely (SR5).
For two unbounded coordinates use the complete polynomial identity
(WR6), replace $X,Y$ by their normalized expressions, and keep
$P=pZ$, $Q=qZ$ as separate roots. The full residue on
$(W,\widehat X,\widehat Y,P,Q)$ is diagonal with pivots
\[
 (C_0,1,1,D_0,m_0+1/2)>0.                         \tag{SR8}
\]
All exact coefficients are bounded. The finite lift applies in
$p,q$, also at their literal zero values. At $p=q=0$ both
$P,Q$ vanish, preserving the original exact null state direction.

Every bounded pair has both residues zero, $U>0$, or $U=0,V>0$.
These are covered by the local lemma, (SR3), and (SR3) or (SR4).
Single infinities use (SR5)--(SR7); double infinity uses (SR8).
Restoration away from a literal endpoint multiplies by a pure-testing
square; no inverse is taken at the endpoint. Every state coordinate
and first normal was retained. Applying the real identity separately
to real and imaginary state parts gives the complex-state result.
No nonemptiness of a residual intersection is asserted.

\begin{corollary}[Two source-unit assumptions removed]
In the revision-9 weak-C2 source theorem, both additional requirements
that $r$ and $z/B$ be bounded units can be omitted, provided $1-r$
and the other original gaps and ratios remain positive bounded units.
\end{corollary}
\begin{proof}
The original $0<r<1$ is either a unit or infinitesimal.
For unit $r$, revisions 9 and 10 cover unit and infinitesimal $z/B$.
For infinitesimal $r$, the present theorem covers both possibilities.
Positivity of $G/B=1-(z/B)(1-r-H)$ and the retained unit $1-r$
exclude unbounded $z/B$. These source-only alternatives are exhaustive.
\end{proof}
The remaining $1-r$-infinitesimal and other source-gap openings
are not admitted by this corollary. The failed simple rescaling on
that opposite boundary is recorded separately.
